
\documentclass{article}
\usepackage{colortbl}
\usepackage{makecell}
\usepackage{multirow}
\usepackage{supertabular}

\begin{document}

\newcounter{utterance}

\centering \large Interaction Transcript for game `dond', experiment `coop\_en', episode 14 with Qwen3{-}30B{-}A3B{-}t0.0.
\vspace{24pt}

{ \footnotesize  \setcounter{utterance}{1}
\setlength{\tabcolsep}{0pt}
\begin{supertabular}{c@{$\;$}|p{.15\linewidth}@{}p{.15\linewidth}p{.15\linewidth}p{.15\linewidth}p{.15\linewidth}p{.15\linewidth}}
    \# & $\;$A & \multicolumn{4}{c}{Game Master} & $\;\:$B\\
    \hline

    \theutterance \stepcounter{utterance}  
    & & \multicolumn{4}{p{0.6\linewidth}}{
        \cellcolor[rgb]{0.9,0.9,0.9}{
            \makecell[{{p{\linewidth}}}]{
                \texttt{\tiny{[P1$\langle$GM]}}
                \texttt{You are playing a cooperative negotiation game in which you have to agree on how to divide a set of items among you and another player.} \\
\\ 
\texttt{Rules:} \\
\texttt{(a) You and the other player are given a set of items. Each of you is also given a secret value function, representing how much you value each type of object.} \\
\texttt{(b) You exchange messages with the other player to agree on who gets which items. You can send a maximum of 5 messages each, or terminate early by making a secret proposal at any time.} \\
\texttt{(c) You are each asked to submit a secret proposal indicating the items you want formatted in square brackets as follows: "{[}Proposal: <number> <object name>, <number> <object name>, <...>{]}"} \\
\texttt{(d) If your proposals are complementary, i.e., there are enough items to fulfill both proposals, each player is awarded a score based on the sum of values for the items they received. Otherwise, both of you get zero points.} \\
\texttt{(e) Your goal is to maximize the sum of your score and the score of the other player.} \\
\\ 
\texttt{Let us start.} \\
\\ 
\texttt{The set of available items is:} \\
\\ 
\texttt{2 mirrors, 1 brick, 3 knives, 1 comb} \\
\\ 
\texttt{Your secret value function is:} \\
\\ 
\texttt{mirror: 1, brick: 4, knife: 1, comb: 1} \\
\\ 
\texttt{IMPORTANT: Your messages, unless it is the secret proposal, are directly transmitted to the other player, so do not include any response to the rules or text announcing your message. To make a secret proposal, use the indicated format. Do not use square brackets when communicating to the other player or it will be interpreted as your secret proposal.} \\
            }
        }
    }
    & \\ \\

    \theutterance \stepcounter{utterance}  
    & \multicolumn{4}{p{0.6\linewidth}}{
        \cellcolor[rgb]{0.9,0.9,0.9}{
            \makecell[{{p{\linewidth}}}]{
                \texttt{\tiny{[P1$\rangle$GM]}}
                \texttt{<think>} \\
\texttt{Okay, let's see. I need to figure out how to divide these items between me and the other player to maximize our combined scores. The items are 2 mirrors, 1 brick, 3 knives, and 1 comb. My value function is mirror:1, brick:4, knife:1, comb:1. So I value the brick the most, then the others are lower. The other player's values aren't given, so I have to guess or figure out based on possible strategies.} \\
\\ 
\texttt{First, the goal is to make sure our proposals are complementary. That means if I ask for certain items, the other player's proposal should have the remaining items. But since I don't know their values, I need to make a proposal that's likely to be compatible.} \\
\\ 
\texttt{Since I value the brick the most (4 points), maybe I should try to get the brick. But if the other player also values the brick, we might conflict. Alternatively, if I take the brick, maybe they can take the mirrors or knives. But how to coordinate?} \\
\\ 
\texttt{Let me think about possible allocations. The total items are 2 mirrors, 1 brick, 3 knives, 1 comb. Let's consider that the other player might have different values. Maybe they value mirrors more? Or maybe they value knives? Without knowing, I have to make an assumption.} \\
\\ 
\texttt{If I propose to take the brick and maybe some other items, but the other player might need the brick. Alternatively, if I can get the brick, that's 4 points for me. If I can get the brick and the comb, that's 5. But maybe the other player would prefer the mirrors. Let's think about possible complementary proposals.} \\
\\ 
\texttt{Suppose I propose to take the brick and the comb. That's 4 + 1 = 5. Then the other player would need to take the remaining items: 2 mirrors, 3 knives. But their value for those might be higher. Alternatively, if I take the brick and some knives, but I don't know.} \\
\\ 
\texttt{Alternatively, maybe I should propose to take the brick and the mirrors. But I only value mirrors at 1 each. So 2 mirrors would be 2 points. But the brick is 4. So total 6. But if the other player can take the knives and comb, which they might value more. But again, without knowing their values, it's tricky.} \\
\\ 
\texttt{Another approach: since the} \\
            }
        }
    }
    & & \\ \\

    \theutterance \stepcounter{utterance}  
    & & & \multicolumn{2}{p{0.3\linewidth}}{
        \cellcolor[rgb]{0.9,0.9,0.9}{
            \makecell[{{p{\linewidth}}}]{
                \texttt{\tiny{[GM$|$GM]}}
                \texttt{continue} \\
            }
        }
    }
    & & \\ \\

    \theutterance \stepcounter{utterance}  
    & & & \multicolumn{2}{p{0.3\linewidth}}{
        \cellcolor[rgb]{0.9,0.9,0.9}{
            \makecell[{{p{\linewidth}}}]{
                \texttt{\tiny{[GM$|$GM]}}
                \texttt{<think>} \\
\texttt{Okay, let's see. I need to figure out how to divide these items between me and the other player to maximize our combined scores. The items are 2 mirrors, 1 brick, 3 knives, and 1 comb. My value function is mirror:1, brick:4, knife:1, comb:1. So I value the brick the most, then the others are lower. The other player's values aren't given, so I have to guess or figure out based on possible strategies.} \\
\\ 
\texttt{First, the goal is to make sure our proposals are complementary. That means if I ask for certain items, the other player's proposal should have the remaining items. But since I don't know their values, I need to make a proposal that's likely to be compatible.} \\
\\ 
\texttt{Since I value the brick the most (4 points), maybe I should try to get the brick. But if the other player also values the brick, we might conflict. Alternatively, if I take the brick, maybe they can take the mirrors or knives. But how to coordinate?} \\
\\ 
\texttt{Let me think about possible allocations. The total items are 2 mirrors, 1 brick, 3 knives, 1 comb. Let's consider that the other player might have different values. Maybe they value mirrors more? Or maybe they value knives? Without knowing, I have to make an assumption.} \\
\\ 
\texttt{If I propose to take the brick and maybe some other items, but the other player might need the brick. Alternatively, if I can get the brick, that's 4 points for me. If I can get the brick and the comb, that's 5. But maybe the other player would prefer the mirrors. Let's think about possible complementary proposals.} \\
\\ 
\texttt{Suppose I propose to take the brick and the comb. That's 4 + 1 = 5. Then the other player would need to take the remaining items: 2 mirrors, 3 knives. But their value for those might be higher. Alternatively, if I take the brick and some knives, but I don't know.} \\
\\ 
\texttt{Alternatively, maybe I should propose to take the brick and the mirrors. But I only value mirrors at 1 each. So 2 mirrors would be 2 points. But the brick is 4. So total 6. But if the other player can take the knives and comb, which they might value more. But again, without knowing their values, it's tricky.} \\
\\ 
\texttt{Another approach: since the} \\
            }
        }
    }
    & & \\ \\

    \theutterance \stepcounter{utterance}  
    & & \multicolumn{4}{p{0.6\linewidth}}{
        \cellcolor[rgb]{0.9,0.9,0.9}{
            \makecell[{{p{\linewidth}}}]{
                \texttt{\tiny{[GM$\rangle$P2]}}
                \texttt{You are playing a cooperative negotiation game in which you have to agree on how to divide a set of items among you and another player.} \\
\\ 
\texttt{Rules:} \\
\texttt{(a) You and the other player are given a set of items. Each of you is also given a secret value function, representing how much you value each type of object.} \\
\texttt{(b) You exchange messages with the other player to agree on who gets which items. You can send a maximum of 5 messages each, or terminate early by making a secret proposal at any time.} \\
\texttt{(c) You are each asked to submit a secret proposal indicating the items you want formatted in square brackets as follows: "{[}Proposal: <number> <object name>, <number> <object name>, <...>{]}"} \\
\texttt{(d) If your proposals are complementary, i.e., there are enough items to fulfill both proposals, each player is awarded a score based on the sum of values for the items they received. Otherwise, both of you get zero points.} \\
\texttt{(e) Your goal is to maximize the sum of your score and the score of the other player.} \\
\\ 
\texttt{Let us start.} \\
\\ 
\texttt{The set of available items is:} \\
\\ 
\texttt{2 mirrors, 1 brick, 3 knives, 1 comb} \\
\\ 
\texttt{Your secret value function is:} \\
\\ 
\texttt{mirror: 2, brick: 1, knife: 0, comb: 5} \\
\\ 
\texttt{IMPORTANT: Your messages, unless it is the secret proposal, are directly transmitted to the other player, so do not include any response to the rules or text announcing your message. To make a secret proposal, use the indicated format. Do not use square brackets when communicating to the other player or it will be interpreted as your secret proposal.} \\
\\ 
\\ 
\texttt{<think>} \\
\texttt{Okay, let's see. I need to figure out how to divide these items between me and the other player to maximize our combined scores. The items are 2 mirrors, 1 brick, 3 knives, and 1 comb. My value function is mirror:1, brick:4, knife:1, comb:1. So I value the brick the most, then the others are lower. The other player's values aren't given, so I have to guess or figure out based on possible strategies.} \\
\\ 
\texttt{First, the goal is to make sure our proposals are complementary. That means if I ask for certain items, the other player's proposal should have the remaining items. But since I don't know their values, I need to make a proposal that's likely to be compatible.} \\
\\ 
\texttt{Since I value the brick the most (4 points), maybe I should try to get the brick. But if the other player also values the brick, we might conflict. Alternatively, if I take the brick, maybe they can take the mirrors or knives. But how to coordinate?} \\
\\ 
\texttt{Let me think about possible allocations. The total items are 2 mirrors, 1 brick, 3 knives, 1 comb. Let's consider that the other player might have different values. Maybe they value mirrors more? Or maybe they value knives? Without knowing, I have to make an assumption.} \\
\\ 
\texttt{If I propose to take the brick and maybe some other items, but the other player might need the brick. Alternatively, if I can get the brick, that's 4 points for me. If I can get the brick and the comb, that's 5. But maybe the other player would prefer the mirrors. Let's think about possible complementary proposals.} \\
\\ 
\texttt{Suppose I propose to take the brick and the comb. That's 4 + 1 = 5. Then the other player would need to take the remaining items: 2 mirrors, 3 knives. But their value for those might be higher. Alternatively, if I take the brick and some knives, but I don't know.} \\
\\ 
\texttt{Alternatively, maybe I should propose to take the brick and the mirrors. But I only value mirrors at 1 each. So 2 mirrors would be 2 points. But the brick is 4. So total 6. But if the other player can take the knives and comb, which they might value more. But again, without knowing their values, it's tricky.} \\
\\ 
\texttt{Another approach: since the} \\
            }
        }
    }
    & \\ \\

    \theutterance \stepcounter{utterance}  
    & & & \multicolumn{4}{p{0.6\linewidth}}{
        \cellcolor[rgb]{0.9,0.9,0.9}{
            \makecell[{{p{\linewidth}}}]{
                \texttt{\tiny{[GM$\langle$P2]}}
                \texttt{<think>} \\
\texttt{Okay, let's try to figure out the best way to divide the items. The key is to make a proposal that's likely to be complementary with the other player's. Since I don't know their values, I need to make an assumption based on my own values and what might be fair.} \\
\\ 
\texttt{My values are: brick (4), mirror (1), knife (1), comb (1). So the brick is the most valuable to me. The other player might have different values, but maybe they also value the brick. If I take the brick, they might need to take other items. Alternatively, if I don't take the brick, maybe they can take it, but I need to ensure that our proposals don't conflict.} \\
\\ 
\texttt{Let me think about possible allocations. If I take the brick and the comb (total 4 + 1 = 5), then the other player would get the mirrors and knives. But if they value mirrors more, that could work. Alternatively, if I take the brick and some knives, but I don't know.} \\
\\ 
\texttt{Wait, the other player's values aren't given. So maybe I should propose something that leaves the most valuable items for them if they have higher values. But how?} \\
\\ 
\texttt{Alternatively, since the total value for me is 2 mirrors (2), 1 brick (4), 3 knives (3), 1 comb (1) = total 10. But the other player's total isn't known. To maximize the sum, we need to allocate items so that both of us get as much as possible.} \\
\\ 
\texttt{But since the proposals need to be complementary, I need to make sure that if I take certain items, the other player can take the rest. So maybe I should propose to take the brick and the comb, which are the highest value items for me. Then the other player would get the mirrors and knives. But if they value mirrors more, that's okay. Alternatively, if they value knives, maybe they can take those.} \\
\\ 
\texttt{Alternatively, maybe I should propose to take the brick and the mirrors. But mirrors are only worth 1 each, so that's 2 points. But the brick is 4, so total 6. Then the other player gets the knives and comb. But the comb is 1 for me, but maybe they value it more.} \\
\\ 
\texttt{But without knowing their values, it's hard. Maybe the safest way is to take the brick, which is the highest value for me} \\
            }
        }
    }
     \\ \\

    \theutterance \stepcounter{utterance}  
    & & & \multicolumn{2}{p{0.3\linewidth}}{
        \cellcolor[rgb]{0.9,0.9,0.9}{
            \makecell[{{p{\linewidth}}}]{
                \texttt{\tiny{[GM$|$GM]}}
                \texttt{continue} \\
            }
        }
    }
    & & \\ \\

    \theutterance \stepcounter{utterance}  
    & & & \multicolumn{2}{p{0.3\linewidth}}{
        \cellcolor[rgb]{0.9,0.9,0.9}{
            \makecell[{{p{\linewidth}}}]{
                \texttt{\tiny{[GM$|$GM]}}
                \texttt{<think>} \\
\texttt{Okay, let's try to figure out the best way to divide the items. The key is to make a proposal that's likely to be complementary with the other player's. Since I don't know their values, I need to make an assumption based on my own values and what might be fair.} \\
\\ 
\texttt{My values are: brick (4), mirror (1), knife (1), comb (1). So the brick is the most valuable to me. The other player might have different values, but maybe they also value the brick. If I take the brick, they might need to take other items. Alternatively, if I don't take the brick, maybe they can take it, but I need to ensure that our proposals don't conflict.} \\
\\ 
\texttt{Let me think about possible allocations. If I take the brick and the comb (total 4 + 1 = 5), then the other player would get the mirrors and knives. But if they value mirrors more, that could work. Alternatively, if I take the brick and some knives, but I don't know.} \\
\\ 
\texttt{Wait, the other player's values aren't given. So maybe I should propose something that leaves the most valuable items for them if they have higher values. But how?} \\
\\ 
\texttt{Alternatively, since the total value for me is 2 mirrors (2), 1 brick (4), 3 knives (3), 1 comb (1) = total 10. But the other player's total isn't known. To maximize the sum, we need to allocate items so that both of us get as much as possible.} \\
\\ 
\texttt{But since the proposals need to be complementary, I need to make sure that if I take certain items, the other player can take the rest. So maybe I should propose to take the brick and the comb, which are the highest value items for me. Then the other player would get the mirrors and knives. But if they value mirrors more, that's okay. Alternatively, if they value knives, maybe they can take those.} \\
\\ 
\texttt{Alternatively, maybe I should propose to take the brick and the mirrors. But mirrors are only worth 1 each, so that's 2 points. But the brick is 4, so total 6. Then the other player gets the knives and comb. But the comb is 1 for me, but maybe they value it more.} \\
\\ 
\texttt{But without knowing their values, it's hard. Maybe the safest way is to take the brick, which is the highest value for me} \\
            }
        }
    }
    & & \\ \\

    \theutterance \stepcounter{utterance}  
    & & \multicolumn{4}{p{0.6\linewidth}}{
        \cellcolor[rgb]{0.9,0.9,0.9}{
            \makecell[{{p{\linewidth}}}]{
                \texttt{\tiny{[P1$\langle$GM]}}
                \texttt{<think>} \\
\texttt{Okay, let's try to figure out the best way to divide the items. The key is to make a proposal that's likely to be complementary with the other player's. Since I don't know their values, I need to make an assumption based on my own values and what might be fair.} \\
\\ 
\texttt{My values are: brick (4), mirror (1), knife (1), comb (1). So the brick is the most valuable to me. The other player might have different values, but maybe they also value the brick. If I take the brick, they might need to take other items. Alternatively, if I don't take the brick, maybe they can take it, but I need to ensure that our proposals don't conflict.} \\
\\ 
\texttt{Let me think about possible allocations. If I take the brick and the comb (total 4 + 1 = 5), then the other player would get the mirrors and knives. But if they value mirrors more, that could work. Alternatively, if I take the brick and some knives, but I don't know.} \\
\\ 
\texttt{Wait, the other player's values aren't given. So maybe I should propose something that leaves the most valuable items for them if they have higher values. But how?} \\
\\ 
\texttt{Alternatively, since the total value for me is 2 mirrors (2), 1 brick (4), 3 knives (3), 1 comb (1) = total 10. But the other player's total isn't known. To maximize the sum, we need to allocate items so that both of us get as much as possible.} \\
\\ 
\texttt{But since the proposals need to be complementary, I need to make sure that if I take certain items, the other player can take the rest. So maybe I should propose to take the brick and the comb, which are the highest value items for me. Then the other player would get the mirrors and knives. But if they value mirrors more, that's okay. Alternatively, if they value knives, maybe they can take those.} \\
\\ 
\texttt{Alternatively, maybe I should propose to take the brick and the mirrors. But mirrors are only worth 1 each, so that's 2 points. But the brick is 4, so total 6. Then the other player gets the knives and comb. But the comb is 1 for me, but maybe they value it more.} \\
\\ 
\texttt{But without knowing their values, it's hard. Maybe the safest way is to take the brick, which is the highest value for me} \\
            }
        }
    }
    & \\ \\

    \theutterance \stepcounter{utterance}  
    & \multicolumn{4}{p{0.6\linewidth}}{
        \cellcolor[rgb]{0.9,0.9,0.9}{
            \makecell[{{p{\linewidth}}}]{
                \texttt{\tiny{[P1$\rangle$GM]}}
                \texttt{<think>} \\
\texttt{Okay, I'll start by proposing to take the brick and the comb, as those are the highest value items for me. This way, I secure 4 + 1 = 5 points. Then, the other player can take the remaining items: 2 mirrors, 3 knives. I hope their values align with that. If not, maybe they'll adjust. But since I can only send 5 messages, I'll go with this proposal.} \\
\\ 
\texttt{{[}Proposal: 1 brick, 1 comb{]}} \\
\texttt{</think>} \\
\\ 
\texttt{{[}Proposal: 1 brick, 1 comb{]}} \\
            }
        }
    }
    & & \\ \\

    \theutterance \stepcounter{utterance}  
    & & & \multicolumn{2}{p{0.3\linewidth}}{
        \cellcolor[rgb]{0.9,0.9,0.9}{
            \makecell[{{p{\linewidth}}}]{
                \texttt{\tiny{[GM$|$GM]}}
                \texttt{proposal} \\
            }
        }
    }
    & & \\ \\

    \theutterance \stepcounter{utterance}  
    & & & \multicolumn{2}{p{0.3\linewidth}}{
        \cellcolor[rgb]{0.9,0.9,0.9}{
            \makecell[{{p{\linewidth}}}]{
                \texttt{\tiny{[GM$|$GM]}}
                [0, 1, 0, 1]
            }
        }
    }
    & & \\ \\

    \theutterance \stepcounter{utterance}  
    & & \multicolumn{4}{p{0.6\linewidth}}{
        \cellcolor[rgb]{0.9,0.9,0.9}{
            \makecell[{{p{\linewidth}}}]{
                \texttt{\tiny{[GM$\rangle$P2]}}
                \texttt{The other player has submitted a secret proposal. It is now your turn to submit a proposal of your own.} \\
            }
        }
    }
    & \\ \\

    \theutterance \stepcounter{utterance}  
    & & & \multicolumn{4}{p{0.6\linewidth}}{
        \cellcolor[rgb]{0.9,0.9,0.9}{
            \makecell[{{p{\linewidth}}}]{
                \texttt{\tiny{[GM$\langle$P2]}}
                \texttt{<think>} \\
\texttt{Okay, the other player has already submitted their secret proposal. Now it's my turn to submit mine. Let me think about what they might have proposed.} \\
\\ 
\texttt{Since I don't know their values, but I know my own, I need to make a proposal that complements theirs. Let's recall the items: 2 mirrors, 1 brick, 3 knives, 1 comb. My values are mirror:1, brick:4, knife:1, comb:1.} \\
\\ 
\texttt{If the other player's proposal is complementary, then together our items should cover all the available items. So whatever they took, I need to take the rest. But since I don't know their proposal, I have to guess.} \\
\\ 
\texttt{Maybe they proposed something like taking the brick and the comb. If that's the case, then I should take the mirrors and knives. Alternatively, if they took the mirrors, I should take the brick and comb.} \\
\\ 
\texttt{But how to decide? Let's think about possible scenarios. If I propose to take the brick and the comb, that's 4 + 1 = 5 points for me. If they took the mirrors and knives, that's 2*1 + 3*1 = 5 points for them. Total 10. That's good.} \\
\\ 
\texttt{Alternatively, if they took the brick, then I should take the comb and other items. But without knowing, it's tricky.} \\
\\ 
\texttt{Wait, the problem says that if the proposals are complementary, we get points. Otherwise, zero. So I need to make sure that my proposal doesn't conflict with theirs. But since I don't know theirs, I have to make an educated guess.} \\
\\ 
\texttt{Maybe the other player is also trying to maximize the total, so they might have proposed something that leaves the most valuable items for me. But I can't be sure.} \\
\\ 
\texttt{Alternatively, since I value the brick the most, maybe I should propose to take the brick and the comb, and hope that they take the rest. Let's go with that. So my proposal would be {[}Proposal: 1 brick, 1 comb{]}. That way, I get the brick (4) and comb (1), totaling 5. The other player would get the mirrors and knives. If they value those, that's okay. The total combined score would be 5 (mine) plus whatever they get from mirrors and knives. Since mirrors are 1 each and knives 1 each, if they take 2 mirrors and} \\
            }
        }
    }
     \\ \\

    \theutterance \stepcounter{utterance}  
    & & & \multicolumn{2}{p{0.3\linewidth}}{
        \cellcolor[rgb]{0.9,0.9,0.9}{
            \makecell[{{p{\linewidth}}}]{
                \texttt{\tiny{[GM$|$GM]}}
                \texttt{proposal} \\
            }
        }
    }
    & & \\ \\

    \theutterance \stepcounter{utterance}  
    & & & \multicolumn{2}{p{0.3\linewidth}}{
        \cellcolor[rgb]{0.9,0.9,0.9}{
            \makecell[{{p{\linewidth}}}]{
                \texttt{\tiny{[GM$|$GM]}}
                [0, 1, 0, 1]
            }
        }
    }
    & & \\ \\

    \theutterance \stepcounter{utterance}  
    & & & \multicolumn{2}{p{0.3\linewidth}}{
        \cellcolor[rgb]{0.9,0.9,0.9}{
            \makecell[{{p{\linewidth}}}]{
                \texttt{\tiny{[GM$|$GM]}}
                [[0, 1, 0, 1], [0, 1, 0, 1]]
            }
        }
    }
    & & \\ \\

\end{supertabular}
}

\end{document}
